\documentclass[11pt]{article}

\usepackage{amsmath,amssymb,amsthm,mathtools}
\usepackage[margin=1in]{geometry}
\usepackage{microtype}

\newtheorem{theorem}{Theorem}[section]
\newtheorem{lemma}[theorem]{Lemma}
\newtheorem{corollary}[theorem]{Corollary}
\newtheorem{remark}[theorem]{Remark}

\newcommand{\dd}{\,\mathrm d}

\title{Orbit-Balanced Leaves on the Diamond Graph:\\
Fractional Degree Bias Beyond Vertex Transitivity}
\author{}
\date{}

\begin{document}

\maketitle

\begin{abstract}
Let $D=K_4-e$ be the diamond graph.  Its automorphism group has two
vertex orbits of size two: the endpoints of the shared triangle edge and
the two page vertices.  Attach a total of $k$ pendant leaves within each
orbit, with arbitrary distribution inside the orbit.  For $1\leq k\leq2$,
we prove the chromatic-polynomial graphon lower bound throughout
$p\geq1/2$.  Orbit symmetrization makes the four root exponents uniformly
$k/2$; a change of measure, fractional H\"older inequality, the accepted
pure-chordal diamond bound, and an elementary scalar monotonicity then give
the exact target.  The case $k=1$ proves the formerly open NetworkX Atlas
graph $113$.
\end{abstract}


%%%%%%%%%%%%%%%%%%%%%%%%%%%%%%%%%%%%%%%%%%%%%%%%%%%%%%%%%%%%%%%%%%%%%%%%%%%%%%%
\section{Family, orbit structure, and target}
%%%%%%%%%%%%%%%%%%%%%%%%%%%%%%%%%%%%%%%%%%%%%%%%%%%%%%%%%%%%%%%%%%%%%%%%%%%%%%%

Let $(\Omega,\mu)$ be an arbitrary probability space and let
$W\colon\Omega^2\to[0,1]$ be symmetric and measurable.  Put
\[
 p=\int_{\Omega^2}W(x,y)\dd\mu(x)\dd\mu(y),
 \qquad
 d(x)=\int_\Omega W(x,y)\dd\mu(y).
\]

Label the diamond $D$ by spine vertices $a,b$ and page vertices $c,d$,
with edge set
\[
 ab,\ ac,\ ad,\ bc,\ bd.
\]
Thus $D$ is the union of the triangles $abc$ and $abd$ along the spine
edge $ab$.  Its automorphism group contains the independent swaps
$a\leftrightarrow b$ and $c\leftrightarrow d$; these are exactly its two
vertex orbits.

Choose nonnegative integers
\[
 s_a+s_b=k,
 \qquad
 q_c+q_d=k,
 \qquad 1\leq k\leq2,
\]
and attach the indicated number of private pendant leaves to each core
vertex.  Denote the resulting graph by
\[
 D_k(s_a,s_b;q_c,q_d).
\]
There are $2k$ leaves in total.  Since
\[
 \chi_D(x)=x(x-1)(x-2)^2,
\]
independent colouring of the leaves gives
\begin{equation}
 \chi_{D_k}(x)=x(x-1)^{2k+1}(x-2)^2,
 \qquad \chi(D_k)=3.
 \label{eq:chromatic}
\end{equation}
Consequently
\begin{equation}
 \Phi_{D_k}(p)=p^{2k+1}(2p-1)^2.
 \label{eq:target}
\end{equation}
This is a polynomial identity, including the continued value one at $p=1$;
at the required threshold $p=1/2$, the target is zero.

We use the accepted pure-chordal theorem only in its diamond instance:
every graphon $V$ of edge density $s\geq1/2$ satisfies
\begin{equation}
 t(D,V)\geq\phi(s),
 \qquad
 \phi(s):=s(2s-1)^2.
 \label{eq:diamond-input}
\end{equation}
Indeed, the two maximal cliques of $D$ are triangles, so $D$ is
$3$-pure chordal.  The theorem applies on arbitrary probability spaces.


%%%%%%%%%%%%%%%%%%%%%%%%%%%%%%%%%%%%%%%%%%%%%%%%%%%%%%%%%%%%%%%%%%%%%%%%%%%%%%%
\section{Fractional edge and scalar transfer}
%%%%%%%%%%%%%%%%%%%%%%%%%%%%%%%%%%%%%%%%%%%%%%%%%%%%%%%%%%%%%%%%%%%%%%%%%%%%%%%

\begin{lemma}[Fractionally degree-weighted edge]
\label{lem:fractional-edge}
For $0<\alpha\leq1$, set
\[
 N_\alpha=\int_{\Omega^2}W(x,y)d(x)^\alpha d(y)^\alpha
 \dd\mu(x)\dd\mu(y).
\]
Then
\[
 N_\alpha\geq p^{1+2\alpha}.
\]
\end{lemma}

\begin{proof}
Define $W(x,y)/d(x)$ to be zero on $\{d(x)=0\}$; Fubini implies that
$W(x,y)$ vanishes for almost every $y$ there.  With
\[
 \theta=\frac1{1+2\alpha},
 \qquad
 \eta=\frac\alpha{1+2\alpha},
\]
one has $\theta+2\eta=1$ and
\[
 W
 =\bigl(Wd(x)^\alpha d(y)^\alpha\bigr)^\theta
  \left(\frac W{d(x)}\right)^\eta
  \left(\frac W{d(y)}\right)^\eta
\]
almost everywhere on the support of $W$.  Weighted H\"older yields
\[
 p\leq N_\alpha^\theta
 \left(\int\frac W{d(x)}\dd\mu^2\right)^\eta
 \left(\int\frac W{d(y)}\dd\mu^2\right)^\eta.
\]
The last two integrals are at most one, so $p\leq N_\alpha^\theta$.
\end{proof}

\begin{lemma}[Diamond rescaling]
\label{lem:scalar}
Let $k\in\{1,2\}$, let $p\in(1/2,1]$, and suppose
$0<z\leq p^{2k}$.  Put
\[
 s_0=p\left(\frac{p^{2k}}z\right)^{1/2}.
\]
If $s_0\leq1$, then
\[
 z\phi(s_0)\geq p^{2k}\phi(p).
\]
\end{lemma}

\begin{proof}
For $s\in[1/2,1]$, put
\[
 \psi(s)=\frac{\phi(s)}{s^2}
 =\frac{(2s-1)^2}s=4s-4+\frac1s.
\]
Then
\[
 \psi'(s)=4-\frac1{s^2}\geq0.
\]
The definition of $s_0$ is equivalent to
\[
 z=p^{2k}\left(\frac p{s_0}\right)^2.
\]
Since $s_0\geq p$, monotonicity of $\psi$ gives
\[
 z\phi(s_0)
 =p^{2k+2}\psi(s_0)
 \geq p^{2k+2}\psi(p)
 =p^{2k}\phi(p).
\]
\end{proof}

The function $\phi$ is nondecreasing on $[1/2,1]$, because
\[
 \phi'(s)=(2s-1)(6s-1)\geq0.
\]


%%%%%%%%%%%%%%%%%%%%%%%%%%%%%%%%%%%%%%%%%%%%%%%%%%%%%%%%%%%%%%%%%%%%%%%%%%%%%%%
\section{The orbit-balanced theorem}
%%%%%%%%%%%%%%%%%%%%%%%%%%%%%%%%%%%%%%%%%%%%%%%%%%%%%%%%%%%%%%%%%%%%%%%%%%%%%%%

\begin{theorem}[Orbit-balanced diamond leaves]
\label{thm:main}
For every $k\in\{1,2\}$, every admissible distribution
$s_a+s_b=q_c+q_d=k$, and every graphon $W$ of edge density
$p\in[1/2,1]$,
\[
 \boxed{
 t\bigl(D_k(s_a,s_b;q_c,q_d),W\bigr)
 \geq p^{2k+1}(2p-1)^2
 =\Phi_{D_k}(p).
 }
\]
\end{theorem}

\begin{proof}
At $p=1/2$ the target is zero.  At $p=1$, one has $W=1$ almost everywhere
and both sides equal one.  Hence assume $p\in(1/2,1)$.

For core variables $x_a,x_b,x_c,x_d$, write
\[
 F_D=W(x_a,x_b)W(x_a,x_c)W(x_a,x_d)
     W(x_b,x_c)W(x_b,x_d).
\]
After integrating the leaves, the density is
\[
 \int F_D
 d(x_a)^{s_a}d(x_b)^{s_b}
 d(x_c)^{q_c}d(x_d)^{q_d}\dd\mu^4.
\]
The independent swaps within the spine and page orbits give four equal
integrals.  Average them and apply the arithmetic--geometric mean inequality
pointwise.  Each core variable then receives the average exponent $k/2$:
\begin{equation}
 t(D_k,W)
 \geq\int F_D\prod_{v\in\{a,b,c,d\}}d(x_v)^{k/2}\dd\mu^4.
 \label{eq:symmetrization}
\end{equation}

Set
\[
 \alpha=\frac k2\in(0,1],
 \qquad
 M=\int_\Omega d(x)^\alpha\dd\mu(x),
 \qquad
 z=M^4.
\]
Since $p>0$, one has $M>0$.  Define the probability measure
\[
 \dd\nu(x)=\frac{d(x)^\alpha}{M}\dd\mu(x).
\]
The right side of \eqref{eq:symmetrization} is
\[
 z\,t(D,W;\nu).
\]
The edge density of $W$ on $(\Omega,\nu)$ is
\[
 s=\frac{N_\alpha}{M^2}.
\]

Concavity of $u^\alpha$ gives
\[
 M\leq p^\alpha,
 \qquad 0<z\leq p^{4\alpha}=p^{2k}.
\]
Lemma~\ref{lem:fractional-edge} gives
\[
 s\geq\frac{p^{1+k}}{M^2}
 =p\left(\frac{p^{2k}}z\right)^{1/2}
 =:s_0.
\]
Because $W\leq1$ and $\nu$ is a probability measure, $s\leq1$, hence
$s_0\leq1$.  Also $s\geq s_0\geq p>1/2$.

Apply the accepted diamond bound \eqref{eq:diamond-input} on the biased
space, monotonicity of $\phi$, and Lemma~\ref{lem:scalar}:
\begin{align*}
 t(D_k,W)
 &\geq z\,t(D,W;\nu)\\
 &\geq z\phi(s)
 \geq z\phi(s_0)
 \geq p^{2k}\phi(p)
 =p^{2k+1}(2p-1)^2.
\end{align*}
This equals \eqref{eq:target}.
\end{proof}

At every balanced complete $\ell$-partite graphon, $\ell\geq2$, the degree
is constant, the orbit symmetrization and measure-change inequalities are
equalities, and the accepted diamond bound is sharp.  Thus the theorem has
all expected chromatic equality examples, including the zero target at
$\ell=2$.


%%%%%%%%%%%%%%%%%%%%%%%%%%%%%%%%%%%%%%%%%%%%%%%%%%%%%%%%%%%%%%%%%%%%%%%%%%%%%%%
\section{Atlas 113, formalization audit, and limitations}
%%%%%%%%%%%%%%%%%%%%%%%%%%%%%%%%%%%%%%%%%%%%%%%%%%%%%%%%%%%%%%%%%%%%%%%%%%%%%%%

\begin{corollary}[Atlas 113]
NetworkX Atlas graph $113$, with graph6 code \texttt{E\char94\_O}, is
positive.  It has order six, size seven, chromatic number three, and
\[
 \chi_{F_{113}}(x)=x(x-1)^3(x-2)^2.
\]
Every graphon of edge density $p\in[1/2,1]$ satisfies
\[
 t(F_{113},W)\geq p^3(2p-1)^2=\Phi_{F_{113}}(p).
\]
\end{corollary}

\begin{proof}
Its NetworkX edge list is
\[
 02,\ 03,\ 04,\ 12,\ 13,\ 23,\ 35.
\]
Vertices $0,1,2,3$ induce the diamond with spine $2,3$ and pages $0,1$.
Vertex $4$ is a leaf at page $0$, while vertex $5$ is a leaf at spine
vertex $3$.  Hence this is the case $k=1$ of Theorem~\ref{thm:main}.
\end{proof}

All integrands are bounded, nonnegative, and measurable, so Tonelli--Fubini
applies directly.  The quotient on $\{d=0\}$ and positivity of the biased
measure's normalizing constant are handled explicitly.  No regularity,
minimum-degree assumption, finite-step reduction, injective-copy
interpretation, or numerical estimate is used.  The only accepted graphon
input is the pure-chordal diamond bound with exactly matching hypotheses.

The orbit-balance condition is essential to this proof: the total number of
leaves on the spine orbit must equal the total on the page orbit, so that all
four core variables acquire one common exponent.  The restriction $k\leq2$
is exactly what makes that exponent at most one and permits the concavity
step.  The theorem does not cover a leaf distribution concentrated in only
one orbit, a two-edge tail, or attachments to an arbitrary pure-chordal core.

The finite identities, symmetrized exponents, scalar derivative, and Atlas
identification are reconstructed exactly by
\begin{verbatim}
python codes/verify_diamond_orbit_balanced_leaves.py
\end{verbatim}
For formal verification, the new reusable inputs are the two-orbit
arithmetic--geometric symmetrization and the scalar transfer for
$s(2s-1)^2/s^2$.  The graph-specific module awaiting inspection is
\texttt{lean/Taeyoung/Examples/Graph113.lean}.


%%%%%%%%%%%%%%%%%%%%%%%%%%%%%%%%%%%%%%%%%%%%%%%%%%%%%%%%%%%%%%%%%%%%%%%%%%%%%%%
\section{What the Lean formalization actually proves}
%%%%%%%%%%%%%%%%%%%%%%%%%%%%%%%%%%%%%%%%%%%%%%%%%%%%%%%%%%%%%%%%%%%%%%%%%%%%%%%

The machine-checked development is
\texttt{lean/Taeyoung/Methods/CliqueDist/Diamond.lean}, ending at
\texttt{satisfiesLowerBound\_113}.  It proves the case $k=1$ and only that
case, and it departs from this note at the same two points as the companion
\texttt{notes/clique\_distributed\_leaves.tex} does.

\paragraph{Scope: $k=1$.}
The theorem here covers $1\leq k\leq2$.  The formalization fixes $k=1$, where
the symmetrized exponent $k/2$ is $1/2$ and the biased density is $\sqrt d$; at
$k=2$ it would be $1$, a different bias, and no six-vertex Atlas row needs it.

\paragraph{The orbit average becomes a single exchange.}
This note averages over both orbits, producing four monomials
$d_{a}d_{b}$ with $a\in\{0,1\}$, $b\in\{2,3\}$, and applies AM--GM to all four.
The formalization uses one: swapping $0\leftrightarrow1$ and
$2\leftrightarrow3$ simultaneously is an automorphism of $D$, so the graph with
leaves at $\{0,2\}$ is isomorphic to the one with leaves at $\{1,3\}$, and
\[
 2\,t(H,W)=\int F\bigl(d_0d_2+d_1d_3\bigr)
 \geq\int F\cdot2\sqrt{d_0d_1d_2d_3},
\]
the two-term AM--GM.  The four-term average and the two-term one give the same
weight $\prod_i d_i^{1/2}$; the two-term version needs one graph isomorphism
(\texttt{CliqueDist.isoOrbit}, checked by \texttt{decide}) instead of four
peelings.

\paragraph{The scalar transfer becomes a factorization.}
Lemma~\ref{lem:scalar} is proved here by differentiating $A_D(s)/s^{2}$.  The
formalization instead uses, with $m=M^2$ and $s$ the biased edge density,
\begin{equation}
 p(2s-1)^2-s(2p-1)^2=(s-p)(4ps-1),
 \label{eq:lean-diamond-rescale}
\end{equation}
a \texttt{ring} identity whose two factors are nonnegative because
$s\geq p\geq1/2$.  Together with $m^2s^2\geq p^4$ this gives
$m^2A_D(s)\geq p^2A_D(p)$ directly.

\paragraph{Everything else is reused verbatim.}
Lemma~\ref{lem:fractional-edge} at $\alpha=1/2$ is
\texttt{GeometricMean.sq\_le\_integral\_geoIntegrand}; $M\leq p^{1/2}$ is
$M^2\leq p$, one Cauchy--Schwarz against the constant $1$; the change of measure
is \texttt{CliqueDist.integral\_sqrtDegree\_prod}.  The accepted diamond input
is not re-derived: it is \texttt{BaseCone.base\_diamond}, already verified for
\texttt{notes/six\_verified\_base\_cones.tex}, which is exactly
$t(D,V)\geq z(2z-1)^2$ on $z\geq1/2$.

\end{document}
